\rtmtight
\begin{tblr}{width=\textwidth,colspec={|Q[c,m,2.1cm]|X[l,m]|},hlines,vlines,hspan=minimal,rowsep=1.5pt,colsep=3pt,row{1}={bg=headgray,font=\bfseries}}
\SetCell{c,m}\hfil\logo\hfil & \textbf{POLITEKNIK NEGERI MALANG}\par\textbf{JURUSAN TEKNOLOGI INFORMASI}\par\textbf{PROGRAM STUDI: D4 TEKNIK INFORMATIKA} \\
\SetCell[c=2]{l} \textbf{RENCANA TUGAS MAHASISWA (RTM) DAN RUBRIK PENILAIAN} & \\
Nama Mata Kuliah & Statistik Komputasi \\
Kode Mata Kuliah & RTI254008 \quad \textbf{sks / Jam perminggu:} 2 SKS/4 jam \quad \textbf{Semester:} 4 \\
Dosen Pengampu & Tim Pengajar Program Studi D4 TI \\
\end{tblr}
\assessmentblock{Tugas Analisis Statistik}{Case Method dan praktik komputasi}{SCPMK0707-03501, SCPMK0902-03502}{Mahasiswa menyelesaikan tugas analisis statistik menggunakan Python/R sesuai Sub-CPMK dan konteks mata kuliah.}{Menganalisis dataset, mengimplementasikan metode statistik, menguji hasil, mendokumentasikan kode dan evidence, serta mempresentasikan temuan.}{Kode Python/R, laporan analisis, visualisasi, dan/atau bahan presentasi.}{Ketepatan implementasi metode statistik & 12 & Metode statistika deskriptif dan inferensial diimplementasikan dengan benar, menghasilkan output yang valid sesuai kasus analisis. & Metode utama sudah benar, tetapi sebagian parameter atau interpretasi hasil belum tepat. & Metode tidak tepat atau implementasi menghasilkan output yang tidak dapat diinterpretasikan. \\ \\
Analisis dan interpretasi hasil & 11 & Hasil analisis diinterpretasikan secara akurat, dikaitkan dengan konteks masalah, dan didukung visualisasi yang informatif. & Interpretasi sebagian besar tepat, tetapi kaitan dengan konteks masalah atau visualisasi masih kurang. & Hasil tidak diinterpretasikan atau interpretasi tidak terkait dengan konteks analisis. \\ \\
Dokumentasi kode dan laporan & 7 & Kode terdokumentasi, terstruktur, dan laporan analisis menjelaskan metodologi, temuan, serta kesimpulan secara runtut. & Dokumentasi dan laporan tersedia, tetapi metodologi atau kesimpulan masih kurang lengkap. & Kode tidak terdokumentasi atau laporan tidak menunjukkan alur analisis yang jelas. \\}

\assessmentblock{Kuis Statistika dan Komputasi}{Kuis}{SCPMK0902-03502, SCPMK0707-03501}{Mahasiswa menyelesaikan kuis untuk mengukur pemahaman konsep statistika dan kemampuan komputasional.}{Menjawab soal pilihan ganda dan/atau uraian terkait konsep statistika, probabilitas, distribusi, dan implementasi komputasi.}{Jawaban tertulis kuis.}{Penguasaan konsep statistika & 6 & Konsep probabilitas, distribusi, inferensi, dan pengujian hipotesis dijelaskan dengan benar disertai contoh yang tepat. & Sebagian besar konsep dipahami, tetapi penerapan pada contoh atau penjelasan masih kurang lengkap. & Konsep tidak dipahami atau contoh penerapan keliru. \\ \\
Penerapan metode komputasi & 4 & Langkah komputasi diterapkan dengan tepat, output diinterpretasikan benar, dan kesimpulan analitik konsisten. & Langkah komputasi sebagian besar benar, tetapi interpretasi atau kesimpulan masih terbatas. & Langkah komputasi tidak tepat atau tidak menunjukkan penguasaan metode. \\}

\assessmentblock{UTS Probabilitas dan Inferensi Statistik}{UTS berbasis analisis komputasi}{SCPMK0902-03502, SCPMK0707-03501}{Evaluasi tengah semester untuk mengukur kemampuan analisis probabilitas, distribusi, dan inferensi statistik menggunakan Python/R.}{Menyelesaikan soal analisis menggunakan Python/R, menginterpretasikan hasil, dan mendokumentasikan proses komputasi.}{Kode Python/R, output analisis, dan penjelasan interpretasi.}{Ketepatan probabilitas dan distribusi & 6 & Distribusi dipilih dengan tepat, parameter dihitung benar, dan hasil komputasi divalidasi sesuai teori. & Distribusi dan perhitungan sebagian besar benar, tetapi validasi atau ketepatan parameter masih kurang. & Distribusi salah dipilih atau perhitungan menghasilkan nilai yang tidak masuk akal. \\ \\
Implementasi inferensi statistik & 5 & Pengujian hipotesis diimplementasikan dengan benar, kesimpulan statistik sesuai nilai p dan alpha yang ditetapkan. & Pengujian dilakukan, tetapi interpretasi nilai p atau kesimpulan statistik masih keliru. & Pengujian tidak dilakukan atau kesimpulan tidak berdasarkan hasil statistik. \\ \\
Kualitas solusi komputasional & 4 & Kode bersih, terstruktur, dan mendokumentasikan setiap langkah analisis dengan jelas. & Kode berfungsi, tetapi struktur atau dokumentasi langkah analisis masih terbatas. & Kode tidak berjalan atau tidak menunjukkan proses analisis yang sistematis. \\}

\assessmentblock{PBL Analisis Dataset TI}{Project Based Learning}{SCPMK0707-03501}{Mahasiswa melakukan analisis statistik end-to-end pada dataset TI menggunakan Python/R, mencakup EDA, pemodelan, dan interpretasi.}{Memilih dataset, melakukan EDA, menerapkan metode statistik lanjut, menginterpretasikan temuan, dan mempresentasikan hasil analisis.}{Notebook Python/R, laporan analisis lengkap, visualisasi, dan bahan presentasi.}{Kualitas analisis data eksploratori & 8 & EDA komprehensif: distribusi, outlier, korelasi, dan pola data teridentifikasi dengan visualisasi yang informatif. & EDA mencakup analisis utama, tetapi sebagian pola atau visualisasi masih kurang mendalam. & EDA tidak sistematis atau tidak mengungkap pola yang relevan dari dataset. \\ \\
Implementasi metode statistik lanjut & 8 & Regresi, ANOVA, atau simulasi diterapkan dengan tepat, asumsi divalidasi, dan hasil diinterpretasikan secara kontekstual. & Metode diterapkan, tetapi validasi asumsi atau interpretasi kontekstual masih kurang. & Metode tidak tepat atau tidak ada validasi asumsi yang mendukung hasil analisis. \\ \\
Presentasi dan komunikasi hasil & 6 & Temuan analisis dikomunikasikan secara jelas, didukung visualisasi, dan rekomendasi berbasis data disajikan meyakinkan. & Presentasi cukup jelas, tetapi visualisasi atau rekomendasi masih terbatas atau kurang kuat. & Temuan tidak dikomunikasikan secara efektif atau tidak ada rekomendasi berbasis data. \\}

\assessmentblock{UAS Presentasi Hasil Analisis Statistik}{UAS berbasis presentasi dan defense}{SCPMK0707-03501, SCPMK0902-03502}{Mahasiswa mempresentasikan dan mempertahankan hasil proyek analisis statistik komprehensif menggunakan Python/R.}{Menyiapkan presentasi proyek analisis, mendemonstrasikan hasil komputasi, dan mempertahankan keputusan metodologis.}{Slide presentasi, notebook Python/R lengkap, laporan akhir, dan visualisasi hasil.}{Ketepatan metode dan analisis & 9 & Metode statistik dipilih dengan tepat, diimplementasikan benar, dan hasil analisis dapat direproduksi serta divalidasi. & Metode sebagian besar tepat, tetapi reproduktibilitas atau validasi hasil masih terbatas. & Metode tidak tepat atau hasil analisis tidak dapat divalidasi. \\ \\
Interpretasi dan pengambilan keputusan & 8 & Temuan diinterpretasikan secara akurat, implikasi praktis dirumuskan, dan rekomendasi keputusan didukung bukti statistik. & Interpretasi cukup akurat, tetapi implikasi atau rekomendasi masih kurang spesifik. & Temuan tidak diinterpretasikan atau rekomendasi tidak berdasarkan bukti statistik. \\ \\
Kualitas presentasi dan komunikasi & 6 & Presentasi terstruktur, visualisasi informatif, dan mahasiswa mampu menjawab pertanyaan teknis dengan tepat. & Presentasi cukup jelas, tetapi visualisasi atau kemampuan menjawab pertanyaan masih terbatas. & Presentasi tidak terstruktur atau mahasiswa tidak mampu menjelaskan keputusan metodologis. \\}

\begin{tblr}{width=\textwidth,colspec={|Q[l,m,3.2cm]|X[l,m]|},hlines,vlines,hspan=minimal,row{1-3}={bg=headgray},rowsep=1.5pt,colsep=3pt}
Jadwal Pelaksanaan: & Minggu 1--17 sesuai RPS. Setiap evaluasi memiliki RTM dan rubrik multi-kriteria. \\
Lain-Lain Yang Diperlukan: & LMS, repositori/notebook Python/R, perangkat lunak sesuai mata kuliah, dan perangkat presentasi. \\
Pustaka: & Mengacu pustaka utama dan pendukung pada bagian RPS. \\
\end{tblr}
